\insertImage{92b04921-76d0-42f4-a88e-21248e3e9194.png}{Aquarelle d'un globe terrestre avec des points marqués comme de petits bâtiments à différents endroits, symbolisant des entreprises libérées dans le monde entier.}

\part*{Édito}
Chers lecteurs,

Embarquons ensemble dans le dédale fascinant de l'univers entrepreneurial où, parmi les innombrables formes d'organisation, l'entreprise libérée brille comme un phare de promesses. On peut facilement être ébloui par la perspective de travailleurs autonomes, d'initiatives spontanées et d'un environnement professionnel débordant de liberté et de respect. Oui, ça sonne bien sur le papier, n'est-ce pas ?

C'est ce que j'ai longtemps cru et c'est toujours une conviction.

Mais pourquoi ecrire un livre sur les zones d'ombres d'une forme d'organisation auquel je tiens tant?

Au détour de mes discussions, j'ai eu l'occasion de rencontrer des personnes qui détestent littéralement ce mode d'organisation parceque lorsqu'on le voile du "hype" et "l'espoir d'utopiste", ce mode d'organisation révèlerait une bureaucratie qui accentue le surmenage, la discrimination et le "gaslighting" et qu'accessoirement, ce modèle n'est pas capable de passer à l'échelle.

Cette discussion m'a ouvert les yeux, non pas sur les vis cachés d'un modèle, mais sur les erreurs que l'on peut commettre pensant que ce modèle libérateur n'est pas sans risques, et comme toute mise en place d'un nouveau modèle d'organisation, il necessite de la rélféxion, de la vigilance et surtout de l'encadrement pour éviter les échecs, les erreurs et les dérives.

Comme le dit le vieil adage, toute médaille a son revers. Derrière les belles histoires, les discours inspirants et les études de cas prometteuses, se cache une réalité plus complexe et souvent plus sombre. Quelle est la véritable essence de l'entreprise libérée ? Quels sont ses principes fondateurs ? Comment est-elle gérée au quotidien ? Quels sont ses avantages et ses inconvénients ? Comment est-elle perçue par ceux qui y travaillent, ceux qui la dirigent, ceux qui la conseillent ?

Voici l'objectif de ce livre : éclairer les zones d'ombre, démystifier les clichés, comprendre les enjeux et mettre en lumière les difficultés inhérentes au modèle d'entreprise libérée. Car, chers lecteurs, il est important de comprendre que si l'entreprise libérée a le potentiel de transformer positivement la vie de travail, elle n'est pas une panacée universelle. Il n'y a pas de formule magique en matière d'organisation d'entreprise. Chaque modèle a ses forces et ses faiblesses, et il est essentiel d'aborder ces questions avec un esprit ouvert et critique.

Ce voyage vous guidera à travers les coulisses des entreprises libérées, en explorant les réalités du terrain, en analysant les réussites et les échecs, et en dévoilant les vérités qui se cachent derrière le jargon et les slogans. Nous plongerons dans le vif du sujet, en exposant les défis liés à la communication, à la gestion des ressources humaines, à la performance et à la rentabilité. Nous explorerons également les raisons des échecs, les leçons tirées, ainsi que les critiques et les alternatives à ce modèle.

Alors, préparez-vous à un voyage rempli d'histoires captivantes, de faits surprenants, d'analyses perspicaces et de réflexions profondes. Qu'il s'agisse de redéfinir votre façon de penser le travail, de remodeler votre organisation ou simplement de satisfaire votre curiosité, ce livre a quelque chose pour vous. Attachez bien vos ceintures, car nous sommes sur le point de décoller pour une exploration passionnante de l'envers du décor des entreprises libérées.

Alors installez-vous confortablement, et laissez-vous guider. Mais attention, la lecture de ce livre pourrait bien bouleverser vos convictions et changer votre regard sur le monde de l'entreprise. Vous voilà prévenus !
